\documentclass[12pt,a4paper]{article}
\usepackage[margin=25mm, headheight=15pt, headsep=10pt]{geometry}
\usepackage{fontspec}
\usepackage[table]{xcolor} % Note: table option is required for rowcolors
\usepackage{titlesec}
\usepackage{tocloft}
\usepackage{fancyhdr}
\usepackage{booktabs}
\usepackage{tcolorbox}
\tcbuselibrary{skins, breakable}
\usepackage[colorlinks=true, linkcolor=emerald, urlcolor=emerald, bookmarks=true]{hyperref}

\definecolor{emerald}{RGB}{0,102,51}
\definecolor{darkred}{RGB}{139,0,0}
\definecolor{lightgray}{RGB}{245,245,245}

\usepackage[nil]{babel}
\babelprovide[import, main]{bengali}
\babelprovide[import]{arabic}
\babelfont{rm}[Renderer=HarfBuzz,Scale=1.0]{Noto Serif Bengali}
\babelfont{sf}[Renderer=HarfBuzz]{Noto Sans Bengali}
\babelfont{tt}[Renderer=HarfBuzz]{Noto Sans Bengali}
\babelfont[arabic]{rm}[Renderer=HarfBuzz,Scale=1.1]{Noto Naskh Arabic}

% Section styling
\titleformat{\section}{\Large\bfseries\color{emerald}}{\thesection.}{0.5em}{}
\titleformat{\subsection}{\large\bfseries\color{emerald!85!black}}{\thesubsection.}{0.5em}{}
\titleformat{\subsubsection}{\normalsize\bfseries\color{emerald!70!black}}{\thesubsubsection.}{0.5em}{}
\titlespacing*{\section}{0pt}{18pt}{8pt}

% Fancy Header/Footer setup
\pagestyle{fancy}
\fancyhf{} % clear all header and footer fields
\fancyhead[L]{\color{emerald}\small\textbf{\nouppercase{\leftmark}}}
\fancyfoot[L]{\scriptsize\color{black!50}{\lat{AI}}-সংকলিত, আলেম-যাচাইকৃত নয়}
\fancyfoot[C]{\color{emerald!70!black}\thepage}
\renewcommand{\headrulewidth}{0.4pt}
\renewcommand{\headrule}{\hbox to\headwidth{\color{emerald}\leaders\hrule height \headrulewidth\hfill}}
\renewcommand{\footrulewidth}{0pt}

% Custom boxes
% 1. Ayat/Hadith Box
\newtcolorbox{ayatbox}[1][]{
    enhanced, breakable, 
    colback=emerald!5, 
    colframe=emerald, 
    boxrule=0.5pt, 
    arc=3pt, 
    left=10pt, right=10pt, top=10pt, bottom=10pt,
    #1
}

% 2. Quote/Translation Box
\newtcolorbox{quotebox}[1][]{
    enhanced, breakable,
    frame hidden,
    colback=lightgray,
    borderline west={3pt}{0pt}{emerald},
    left=15pt, right=10pt, top=10pt, bottom=10pt,
    sharp corners,
    #1
}

% 3. AI-generated notice box
\newtcolorbox{warnbox}[1][]{
    enhanced, breakable,
    colback=red!4,
    colframe=darkred,
    boxrule=0.8pt,
    arc=3pt,
    left=10pt, right=10pt, top=10pt, bottom=10pt,
    #1
}

% Commands
\newcommand{\arab}[1]{\foreignlanguage{arabic}{#1}}
\newcommand{\kib}[1]{\textcolor{emerald}{\textbf{#1}}}

% Latin (English) fallback: Noto Serif Bengali has NO Latin glyphs
\newfontfamily\latfont[Renderer=HarfBuzz]{Noto Serif}
\newcommand{\lat}[1]{{\latfont #1}}

\setlength{\parskip}{5pt}
\setlength{\parindent}{0pt}

% Fix for missing bullet in Noto Serif Bengali
\renewcommand{\labelitemi}{$\bullet$}



\begin{document}
\begin{center}{\Large\bfseries\color{emerald} পার্ট ৩ — নামাজের কাজগুলোর শ্রেণিবিভাগ: ফরজ, ওয়াজিব, সুন্নত, মুস্তাহাব, মাকরূহ}\end{center}
\vspace{1em}
\section*{পার্ট ৩ — নামাজের কাজগুলোর শ্রেণিবিভাগ: ফরজ, ওয়াজিব, সুন্নত, মুস্তাহাব, মাকরূহ}

\begin{warnbox}
 \textbf{গুরুত্বপূর্ণ নোটিশ (\lat{Important} \lat{Notice}):} এই কন্টেন্টটি \lat{AI} (কৃত্রিম বুদ্ধিমত্তা) দিয়ে প্রস্তুত ও সংকলিত। \lat{AI} কঠোর সূত্র-মান নীতি মেনে শুধু নির্ভরযোগ্য উৎস থেকে তথ্য সংগ্রহ করেছে — কেবল সহীহ/হাসান হাদিস ও সহীহ সনদের আসার রাখা হয়েছে, দুর্বল সনদ বাদ দেওয়া হয়েছে। তবে এটি কোনো আলেম বা মুহাদ্দিস কর্তৃক যাচাইকৃত নয়।
\end{warnbox}

\begin{quote}
\textbf{সম্পর্কিত ফাইল:} পার্ট ২ (পূর্বশর্ত), পার্ট ৪ (পূর্ণ পদ্ধতি), পার্ট ৫খ (সাহু সিজদা), পার্ট ৬ (মাযহাব-তুলনা), পার্ট ৭ক (নামাজ ভঙ্গকারী)।
\end{quote}

\begin{quote}
\textbf{যে প্রশ্নের উত্তর এই অংশে:} নামাজের ভেতরের প্রতিটি কাজের মর্যাদা (\lat{status}) কী — কোনটি ছাড়লে নামাজ বাতিল, কোনটি ছাড়লে সাহু সিজদা, কোনটি ছাড়লে ফজিলত কমে? আর ফকিহরা কেন এ নিয়ে মতভেদ করেছেন?
\end{quote}

\medskip\hrule\medskip

\subsection*{১. কেন এই শ্রেণিবিভাগ শেখা জরুরি (\lat{Why} \lat{It} \lat{Matters})}

নামাজের ভেতরে বহু কাজ আছে — তাকবির, ফাতিহা, রুকু, সিজদা, তাশাহহুদ, দুআ, সালাম... \textbf{এগুলোর সব এক-ই স্তরের নয়।} ফকিহরা পরিভাষা দিয়ে এগুলো চেনান:

\begin{itemize}
\item \textbf{ফরজ (\lat{Fard}):} ছাড়লে \textbf{নামাজ বাতিল} — ইচ্ছাকৃত বা ভুলে।
\item \textbf{ওয়াজিব (\lat{Wajib}):} হানাফি পরিভাষায় — ছাড়লে নামাজ \textbf{বাতিল নয়}, কিন্তু \textbf{সাহু সিজদা} দেবে (ভুলে ছুটলে); ইচ্ছাকৃত ছাড়লে গুনাহ।
\item \textbf{সুন্নত (\lat{Sunnah}):} নবীজির (সা.) নিয়মিত আমল — ছাড়লে নামাজ বৈধ, কিন্তু বড় ক্ষতি — পূর্ণতার সুবর্ণ সুযোগ হাতছাড়া।
\item \textbf{মুস্তাহাব (\lat{Mustahab}):} উত্তম; না করলে তেমন কিছু যায়-আসে না।
\item \textbf{মাকরূহ (\lat{Makruh}):} অপছন্দনীয় — কমিয়ে দেয় নামাজের সওয়াব ও মনোযোগ।
\item \textbf{মুফসিদ/হারাম:} নামাজ \textbf{ভেঙে দেয়} (পার্ট ৭ক)।
\end{itemize}
\begin{quote}
\textbf{মারাত্মক ভুল ধারণা বাদ দিন:} কেউ কেউ ভাবেন "ফজরের চেয়ে ওয়াজিব নিচু, তাই হয়তো পড়লে হলো"। বাস্তবে \textbf{ওয়াজিব হানাফি মাযহাবে অত্যন্ত গুরুত্বপূর্ণ} — ইচ্ছাকৃত ছাড়লে নামাজ দুর্বল ও গুনাহ হয়; ভুলে ছুটলে সাহু সিজদা লাগে। শুধু ফরজের মতো তা \textbf{একেবারে বাতিল} করে না।
\end{quote}

\medskip\hrule\medskip

\subsection*{২. পরিভাষার রহস্য: কেন চার মাযহাবে "ফরজ/ওয়াজিব" তালিকা অমিল?}

এখানে মূল রহস্যটা \textbf{দলিলের স্তর (\lat{gradation} \lat{of} \lat{evidence})} সম্পর্কে।

\textbf{মূলনীতি:} ফকিহরা দলিলের শক্তি অনুযায়ী হুকুম পাতান। হানাফি মাযহাবে:
\begin{itemize}
\item \textbf{ফরজ} = কুরআনের অকাট্য (কাটিঈ দলালাহ) বা মুতাওয়াতির/মশহুর হাদিসে প্রমাণিত কাজ — যেমন নামাজ নিজেই, রুকু, সিজদা (কারণ এগুলোতে দলিল একেবারে নিশ্চিত)।
\item \textbf{ওয়াজিব} = \textbf{কেবল একক-সনদ (আহাদ) সহীহ হাদিসে} প্রমাণিত কাজ — যেমন ফাতিহা (নিশ্চিত মানে প্রমাণিত), সূরার তিলাওয়াত, তাশাহহুদ, সালাম। এতে \textbf{সন্দেহের (\lat{zann}) মাত্রা} থাকায় হুকুমও ফরজের নিচে।
\end{itemize}
\textbf{জমহুর মাযহাব (শাফেয়ি, মালেকি, হাম্বলি):} তারা \textbf{"ফরজ-ওয়াজিব = একই"} বলেন — কারণ শরিয়তের দলিল যেখানে আসে, সেখানে পালন আবশ্যক, দলিল কাটিঈ না আহাদ হলে তা আলাদা হুকুম তৈরি করে না (তাদের মতে)। তাই জমহুরের কাছে \textbf{ফাতিহা, তাশাহহুদ, সালাম — সব ফরজ}।

\begin{quote}
\textbf{সুতরাং একই কাজ:} হানাফি বলেন "ওয়াজিব", শাফেয়ি বলেন "ফরজ" — পার্থক্য মূলত \textbf{পরিভাষা ও ইত্পত্তির (\lat{deduction}) প্রক্রিয়ায়}, দুই পক্ষই একই দলিলই দেখেন। তবে \textbf{ব্যবহারিক পার্থক্য} হয়ে যায় একটি ক্ষেত্রে — \textbf{ভুলে ছুটে গেলে}: হানাফিতে সাহু সিজদা, জমহুরে (কিছু বিষয়ে) নামাজ পুনরায় বা সাহু সিজদা (বিস্তারিত পার্ট ৫খ)।
\end{quote}

\textbf{এই অধ্যায়ের নকশা:} আমরা কাজগুলোকে \textbf{জমহুরের তালিকা} দিয়ে সাজাব, সাথে প্রতিটির নিচে \textbf{হানাফি অবস্থান} টীকা দেব — যাতে দুই দৃষ্টিকোণই পরিষ্কার দেখা যায়।

\medskip\hrule\medskip

\subsection*{৩. নামাজের রুকন (\lat{Fard}/\lat{Arkan}) — সর্বসম্মত ও মতভেদপূর্ণ}

\subsubsection*{৩.১ সর্বসম্মত রুকনসমূহ (যে ছাড়ালে সুস্পষ্টভাবে ভুল)}

\textbf{সহীহ হাদিস — ভিত্তি (সহীহ বুখারি ৭৫৭; সহীহ মুসলিম ৩৯৭):}

\begin{quote}
আবু হুরাইরা (রাঃ) থেকে: এক ব্যক্তি নামাজ পড়ে এসে বলল, হে আল্লাহর রাসূল! আমার নামাজ সঠিক হলো না। নবীজি (সা.) বললেন: \lat{“}ফিরে গিয়ে আবার নামাজ পড়ো।\lat{”} ... যখন সে আবার পড়ল, নবীজি (সা.) বললেন: \textbf{\lat{“}আমি তোমার মধ্যে রুকনসমূহ ব্যতীত কিছু ত্রুটি দেখছি\lat{”}} (অর্থাৎ রুকনগুলোই মূল; নবীজি (সা.) তাকে প্রতিটি রুকনে স্থিরতা (তুমাঈনাহ) ও সঠিকতা শেখালেন)।
\end{quote}

\textbf{সাধারণ সারণি (জমহুরের তালিকা — রুকন):}

\begin{table}[h]\centering\small
\begin{tabular}{lll}
আন্দোলন/কথা & বিবরণ & হানাফি মত \\
\hline
নিয়ত & হৃদয়ে (পার্ট ২) & ফরজ \\
তাকবিরে তাহরিমা & "আল্লাহু আকবার" & ফরজ \\
কিয়াম (দাঁড়ানো) & কুদরত থাকলে & ফরজ \\
কিরাআত & ফাতিহা + তিলাওয়াত & ফাতিহা = ওয়াজিব; অন্তত একটি আয়াত ফরজ \\
রুকু & ঝুঁকে আল্লাহর মহিমা & ফরজ \\
রুকু থেকে দাঁড়ানো & কিরা "সামিআল্লাহু..." & ফরজ (জমহুর); হানাফিতে ওয়াজিব \\
সিজদা (দুই) & দুই সিজদা + মাঝে বসা & ফরজ \\
শেষ বৈঠক & শেষ রাকাতে তাশাহহুদ & তাশাহহুদ = ওয়াজিব \\
সালাম & ডানে-বামে & ওয়াজিব \\
\hline
\end{tabular}
\end{table}

\subsubsection*{৩.২ যে রুকনগুলোতে বাস্তব ইখতিলাফ আছে}

\#\#\#\# ক. কিরাআত — ফাতিহা কি প্রতিটি রাকাতে আবশ্যক?

\textbf{সহীহ হাদিস (সহীহ বুখারি ৭৫২; সহীহ মুসলিম ৩৯৪):}

\begin{quote}
আবু হুরাইরা (রাঃ) থেকে: নবীজি (সা.) বলেন — \textbf{\lat{“}যে নামাজ পড়ল অথচ ফাতিহা পড়ল না, তার নামাজ খিন্নাহ, অপূর্ণ।\lat{”}} ... বললেন, আল্লাহ বলেন: \textbf{\lat{“}আমি নামাজ ভাগ করেছি আমার ও আমার বান্দার মাঝে অর্ধেক অর্ধেক...\lat{”}} (হাদিসে কুদসি)।
\end{quote}

\begin{quote}
\textbf{জমহুর (শাফেয়ি-মালেকি-হাম্বলি):} প্রতিটি রাকাতে \textbf{ফাতিহা ফরজ}; ছাড়লে নামাজ বাতিল (নিরবচ্ছিন্ন মুক্তাদি ব্যতীত — মালেকি মত)। \textbf{হানাফি:} ফাতিহা \textbf{ওয়াজিব}, শর্ত ফরজের নিচে — প্রতিটি রাকাতে ফরজ হলো \textbf{কুরআন থেকে যতটুকুই যোগ্য একটি অংশ}; হানাফিদের অগ্রাধিকার ফাতিহা + সূরা। \\
\textbf{কারণ:} হানাফিতে ফাতিহা-বাছুরের হাদিসটি "খিন্নাহ/নির্দিষ্ট রাকাতে" প্রযোজ্য — কিন্তু তারা ইমামের পেছনের মুক্তাদিকে \textbf{ফাতিহা পাঠে বাধ্য করে না} ("নিরবচ্ছিন্ন শোনা" বিধানে); জমহুরের কাছে প্রতিটি মুসল্লির নিজের রাকাতে ফাতিহা আবশ্যকতাই প্রাধান্য পায়।
\end{quote}

\#\#\#\# খ. শেষ বৈঠকে তাশাহহুদ ও সালাম

\textbf{সহীহ হাদিস — তাশাহহুদের শিখন (সহীহ বুখারি ৬২৬৫; সহীহ মুসলিম ৪০২):}

\begin{quote}
আবদুল্লাহ ইবনে মাসউদ (রাঃ) থেকে: নবীজি (সা.) আমাকে তাশাহহুদ শিখিয়েছেন — \lat{“}আত-তাহিয়্যাতু লিল্লাহি ওয়াস-সালাওয়াতু ওয়াত-তাইয়িবাত...\lat{”}।
\end{quote}

\textbf{সহীহ হাদিস — সালাম (সহীহ মুসলিম ৪১৪):}

\begin{quote}
নবীজি (সা.) সালাম না দেওয়া পর্যন্ত নামাজ শেষ হতো না — ডানে "আসসালামু আলাইকুম ওয়া রাহমাতুল্লাহ" এবং বামেও।
\end{quote}

\begin{quote}
\textbf{হানাফি:} তাশাহহুদ ও সালাম \textbf{ওয়াজিব} (ভুলে ছুটলে সাহু সিজদা); \textbf{জমহুর:} ফরজ (শেষ বৈঠকে তাশাহহুদ + সালাম ছাড়লে নামাজ বাতিল)। এই পার্থক্যই পার্ট ৫খ-এ সাহু সিজদার নিয়মে ব্যাপকভাবে ফুটে ওঠে।
\end{quote}

\#\#\#\# গ. সিজদায় সাত অঙ্গ

\textbf{সহীহ হাদিস (সহীহ বুখারি ৮১২; সহীহ মুসলিম ৪৯০):}

\begin{quote}
ইবনে আব্বাস (রাঃ) থেকে: নবীজি (সা.) বলেন — \textbf{\lat{“}আমাকে আদেশ করা হয়েছে} (নবীগণের উপর ভারী) সাত অঙ্গে সিজদা করতে: (১) কপাল (নাকসহ), (২) দুই হাতের তালু, (৩) দুই হাঁটু, (৪) দুই পায়ের আঙুল। \textbf{এগুলোর কোনো একটি ছাড়া আমি পছন্দ করি না}\lat{”} — সাত অঙ্গের সিজদায় পূর্ণতা ফরজ/ওয়াজিব, এ নিয়ে সব মাযহাব মূলত একমত; পার্থক্য কেবল \textbf{কপালের নাকের অংশ} ও তালু-স্পর্শের হুকুমে।
\end{quote}

\medskip\hrule\medskip

\subsection*{৪. ওয়াজিব — কোন কাজগুলো হানাফি পরিভাষায় ওয়াজিব (জমহুরে ফরজ)}

নিচের কাজগুলো হানাফি মাযহাব \textbf{ওয়াজিব} বলে; জমহুরের কাছে এগুলোর একই হুকুম। \textbf{ভুলে এগুলো ছুটে গেলে সাহু সিজদা দিতে হয় (হানাফি)} — এটাই প্রধান ব্যবহারিক শর্ত:

\begin{enumerate}
\item \textbf{ফাতিহা} (প্রথম দুই রাকাতে) ও সাথে \textbf{সূরার তিলাওয়াত}।
\item \textbf{রুকু-সিজদায় তাসবিহ} (মূলত ১ বার; মুস্তাহাব)।
\item \textbf{তাদিলে আরকান} — প্রতিটি রুকনে \textbf{স্থিরতা} (তুমাঈনাহ)। [অ-হানাফিদের মধ্যে তুমাঈনাহ সর্বসম্মত ফরজ]
\item \textbf{প্রথম বৈঠক} (২ রাকাত পূর্ণ ৩/৪ রাকাতে)।
\item \textbf{শেষ বৈঠকে তাশাহহুদ}।
\item \textbf{সালাম} (হানাফি: ওয়াজিব; বিতর্কিত ফরজ মতও আছে)।
\item \textbf{সামিআল্লাহু লিমান হামিদাহ / রাব্বানা লাকাল হামদ} — রুকু থেকে উঠার না-কথা।
\item \textbf{বিতরের কুনুত} (বিতর নামাজে — পার্ট ৫ঘ)।
\end{enumerate}
\begin{quote}
\textbf{জমহুরের কাছে উপরের অনেকগুলিই ফরজ} — তাই পার্থক্যটি মূলত \textbf{শত্রু নয়, দলিলের তীব্রতার মাপে}। হানাফি মসজিদে ইমাম ভুলে যেমন "ফাতিহা-শেষে" না বসলে বুঝে নিবেন — সাহু সিজদা দিয়ে ইখতিলাফপূর্ণ জায়গায় সমাধান।
\end{quote}

\medskip\hrule\medskip

\subsection*{৫. সুন্নত ও মুস্তাহাব}

সুন্নত দুই প্রকার:

\begin{itemize}
\item \textbf{সুন্নতে মুআক্কাদা (যুক্ত: কেউ কেউ বলবেন)*- নবীজি (সা.) নিয়মিতভাবে করেন, প্রায় সাথে সাথে ছেড়ে দেননি — যেমন }সুন্নত-রাকাত\textbf{ (যোহরের আগে ৪, ফজরের আগে ২, মাগরিবের পরে ২, ইশার পরে ২ — সংক্ষেপ; বিশদ পার্ট ৫ঘ)। ছাড়লে }নামাজ বৈধ**, তবে বড় ফজিলত হাতছাড়া।
\item \textbf{সুন্নতে গায়রে মুআক্কাদা} — \lat{N-ashab}, তবে নিয়মিত যেমন-তেমন।
\end{itemize}
\textbf{নামাজের ভেতরের গুরুত্বপূর্ণ সুন্নত/মুস্তাহাবসমূহ:}

\begin{enumerate}
\item \textbf{রাফউল ইয়াদাইন} — হাত কান বরাবর তোলা।
\end{enumerate}
\begin{itemize}
\item \textbf{সহীহ হাদিস (সহীহ বুখারি ৭৩৫):} ইবনে উমার (রাঃ) থেকে: নবীজি (সা.) তাকবিরের সময়, রুকুতে যাওয়ার সময় ও রুকু থেকে উঠার সময় \textbf{দুই হাত কাঁধ বরাবর তুলতেন।}
\item \textbf{মাযহাব-অবস্থান:} শাফেয়ি-হাম্বলি: \textbf{প্রতিটি তাকবিরে} (তাকবিরে তাহরিমা, রুকুতে যাওয়া, রুকু থেকে উঠা, তৃতীয় রাকাতে দাঁড়ানো) হাত তোলা \textbf{সুন্নত}। \textbf{হানাফি:} শুধু \textbf{তাকবিরে তাহরিমায়} হাত তোলা সুন্নত; রুকুতে যাওয়ার সময় তোলা \textbf{না} — কারণ আবু হানিফা (রহ.) মনে করতেন তা রহিত (নাসিখ) হয়েছে। \textbf{মালেকি:} শুধু তাকবিরে তাহরিমায় তোলা উত্তম। — \textbf{এটিই সবচেয়ে সুপরিচিত বাস্তব ইখতিলাফ} — কোনো পক্ষ ভুল নয়।
\end{itemize}
\begin{enumerate}
\item \textbf{সানা} (থানা): তাকবিরের পরের দুআ — "সুবহানাকা আল্লাহুম্মা..."।
\item \textbf{আউযুবিল্লাহ-বিসমিল্লাহ} — এবং প্রথম রাকাতে জোরে "বিসমিল্লাহ" (হানাফি-জমহুর-বিচিত্র)।
\item \textbf{আমিন} — ফাতিহা শেষে (জোরে/নিরবে — মাযহাব পার্থক্য)।
\item রুকুতে \textbf{সুবহানা রাব্বিয়াল আযিম} — \textbf{সহীহ হাদিস (সুনানে আবু দাউদ ৮৬৯; সুনানে নাসাঈ ১০৪৫):} \lat{“}সূরায়ে ওয়াকিয়াহ ৭৪ নাযিল হওয়ায় নবীজি (সা.) বললেন — এটাকে তোমাদের রুকুতে বানাও।\lat{”} \textbf{কমপক্ষে ১ বার; মুস্তাহাব ৩ বার।}
\item সিজদায় \textbf{সুবহানা রাব্বিয়াল আলা} — \textbf{সহীহ হাদিস (সুনানে আবু দাউদ ৮৭০):} সূরায়ে আলা ১ নাযিল হলে — \lat{“}তোমাদের সিজদায় বানাও।\lat{”}
\item দুই সিজদার মাঝে \textbf{"রাব্বিগফিরলি"} — (সুনানে আবু দাউদ ৮৭৪-এর মতো বর্ণনায়) মুস্তাহাব।
\item তাশাহহুদের আঙুল ও \textbf{দুরুদ} (সালাতু আলান-নবী) — মুস্তাহাব/সুন্নত।
\item সালামের পরের জিকির — তাসবিহ, তাহমিদ, তাকবির (মুস্তাহাব)।
\end{enumerate}
\begin{quote}
\textbf{মুস্তাহাবের গুরুত্ব:} মুস্তাহাব না করলে গুনাহ নয় — তবে মুস্তাহাবগুলোর একটিও নামাজের মর্যাদা কমায়; মুহাদ্দিস ও ফকিহরা বলেন — একজন মুসল্লি তার সাধ্যমতো \textbf{সব শুদ্ধ-সুন্নত} একত্রে করতে চাইবেন, কেননা সুন্নত-মুস্তাহাব দিয়ে শরিয়তের সৌন্দর্য ও তাসফিয়া (\lat{purification}) বিদ্যায়-আসে।
\end{quote}

\medskip\hrule\medskip

\subsection*{৬. মাকরূহ (অপছন্দনীয়) — যা নামাজের সওয়াব কমায়}

\textbf{মাকরূহ} দুই প্রকার:
\begin{itemize}
\item \textbf{মাকরূহে তাহরিমি (প্রায়-নিষিদ্ধ):} মাযহাবে হারামের সামনে (যেমন ফরজ ছাড়া নিয়মিত \lat{Fil}-এর প্রতি খেয়াল — হানাফিরা বিতর্কিত)। বিশদ পার্ট ৭ক।
\item \textbf{মাকরূহে তানযিহি (হালকা অপছন্দ):} কম ফজিলত — যেমন:
\end{itemize}
\begin{enumerate}
\item নামাজে \textbf{ঘুরে তাকানো} — \textbf{সহীহ হাদিস (সহীহ বুখারি ৭৫০):} \lat{“}লোকেরা চোরা-কাটা করে চুরি... নামাজ থেকে\lat{”} — নবীজি (সা.) বলেছেন এবং স্পষ্টত মানুষ নামাজে এদিক-ওদিক তাকানো অপছন্দ করেন — বরং \lat{“}নামাজে চুরি হলো রুকু-সিজদা না পূর্ণ করা।\lat{”}
\item \textbf{চুল/কাপড়ে হাত বোলানো} (খেঁউ-গুছানো-ব্যবস্থা) — ফকিহদের অপছন্দের তালিকায়।
\item খালি (নিছক) \textbf{পোশাক-ঢিলা নিয়ে মশগুল} হওয়া — যেমন দামি বস্তু-বাদ্য।
\item \textbf{নিজের ও সামনের জায়গায় নকশা/ছবি} থাকা — দৃষ্টি বিক্ষিপ্ত হলে; তবে নাবালেগ-পরিচ্ছন্নতায় শিথিল (সহীহ বুখারি ৪০২-এ আয়েশা (রাঃ)-এর ছবি-চাদরের কাহিনিতে — নবীজি (সা.) সেই চাদরে নামাজ পড়া বন্ধ করেননি, বরং অপছন্দ করেছেন — বিস্তারিত পার্ট ৭ক)।
\item \textbf{পাঞ্জার বিপরীত-পাল্টা} — যেমন ডান পায়ে সিথান, হাতে পানের কণা ইত্যাদি গিলে ফেলা — কিছু মাজহাবে মাকরূহ নয় (চূড়ান্ত মাযহাব-পাথর) — কিন্তু সাধারণত গুরুত্বহীন ঝোঁক।
\item \textbf{বিপদজনক তাড়াহুড়া-চলাচল} — যেমন তাকবিরের সময় ঝুঁকে চলা (সুনানে নাসাঈ — ইমামের পেছনে হুড়মুড়িয়ে) — মাকরূহ।
\end{enumerate}
\begin{quote}
\textbf{মাকরূহ এবং নামাজ বৈধ:} মাকরূহ করলে \textbf{নামাজ বাতিল হয় না}, তবে নামাজের গঠন, পূর্ণতা ও খুশু ব্যাহত হয় — সওয়াব কমে যায়।
\end{quote}

\medskip\hrule\medskip

\subsection*{৭. আসার (সাহাবিদের বর্ণনা) — শ্রেণিকে শিথিলভাবে নেওয়ার শিক্ষা}

\textbf{সহীহ সনদের হাদিস — তুমাঈনাহ'ই রুকন (সুনানে আবু দাউদ ৮৫৬; সহীহ বুখারি ৭৯১):}

\begin{quote}
আবু কাতাদাহ (রাঃ) থেকে: নবীজি (সা.) বলেছেন — \textbf{\lat{“}সে ব্যক্তি মিথ্যা, যে রুকু ও সিজদায় তুমাঈনাহ না পায়... তার নামাজ যথেষ্ট নয় — যতক্ষণ না সে রুকুতে ও সিজদায় (সম্পূর্ণ) স্থির বসে; তারপর (মাথা) উঠায় যতক্ষণ না সোজা দাঁড়িয়ে যায়; তারপর সিজদা করে; তারপর দুই সিজদার মাঝে বৈঠকে স্থির হয়।\lat{”}}
\end{quote}

\begin{quote}
\textbf{শিক্ষা:} ফকিহদের সম্ভাবনার তুলনায় \textbf{নবীজির শিক্ষা সর্বনিম্ন-গুরুত্বহীন চর্চাকে সুসংজ্ঞায়িত করেছে}: শ্রেণিবিভাগ শেখা মানে \textbf{দলিল বুঝা}, একজন সাধারণ মুসল্লিকে যা দরকার — \textbf{তাকবির, কিরাআত, রুকু, সিজদা, স্থিরতা, সালাম} — এগুলো পূর্ণতা দিয়ে পড়া। সুন্নত-মুস্তাহাব দিয়ে পূর্ণতা।
\end{quote}

\medskip\hrule\medskip

\subsection*{৮. চূড়ান্ত চার্ট — এক নজরে সব শ্রেণি}

\begin{table}[h]\centering\small
\begin{tabular}{llll}
শ্রেণি & সংজ্ঞা & ভুলে ছুটে গেলে & ইচ্ছাকৃতভাবে ছাড়লে \\
\hline
\textbf{ফরজ (রুকন)} & অকাট্য দলিলে প্রমাণিত & নামাজ বাতিল — পুনরায় পড়তে হবে & নামাজ বাতিল + গুনাহ \\
\textbf{ওয়াজিব} (হানাফি) & একক-সনদ সহীহ হাদিসে প্রমাণিত & সাহু সিজদা & গুনাহ (সাহু সিজদা নয় — পুনরায় উত্তম) \\
\textbf{সুন্নত} & নবীজির নিয়মিত কর্ম & কিছু যায়-আসে না (নামাজ বৈধ) & ফজিলত হাতছাড়া হয়ে যায় \\
\textbf{মুস্তাহাব} & উত্তম & — & তেমন ক্ষতি নেই \\
\textbf{মাকরূহ} & অপছন্দনীয় & (নামাজ বৈধ) & সওয়াব/খুশু কমে \\
\textbf{মুফসিদ} & নামাজ ভঙ্গকারী & নামাজ ভেঙে যায় — পুনরায় & নামাজ ভেঙে যায় \\
\hline
\end{tabular}
\end{table}

\medskip\hrule\medskip

\subsection*{৯. রেফারেন্স ও উৎসগ্রন্থ}

\begin{itemize}
\item সহীহ বুখারি — ৬৩১, ৭৩৫, ৭৫০, ৭৫২, ৭৫৭, ৭৯১, ৮১২, ৭৫২, ৬২৬৫
\item সহীহ মুসলিম — ৩৯৪, ৩৯৭, ৪০২, ৪১৪, ৪৯০
\item সুনানে আবু দাউদ — ৬১, ৮৫৬, ৮৬৯, ৮৭০, ৮৭৪
\item সুনানে তিরমিজি — ৩ (হাদিস নম্বর — \lat{“}মিফতাহুস-সালাত\lat{”}), ২৩৮
\item সুনানে নাসাঈ — ১০৪৫
\item আল-ফিকহুল ইসলামি ওয়া আদিল্লাতুহু (আল-যুহাইলি) — মাযহাব-তুলনার রেফারেন্স
\item সিফাতু সালাতিন নবী (আলবানি) — রুকন-সুন্নত নির্ধারণের হাদিস বিশ্লেষণ
\end{itemize}
\begin{quote}
\textbf{দ্রষ্টব্য:} এই অংশের মূল লক্ষ্য \textbf{দলিল-ভিত্তিক শ্রেণি-বোধ} তৈরি করা। \lat{“}কোন মাযহাব সঠিক\lat{”} সিদ্ধান্ত দেওয়া নয় — বরং প্রতিটি মাযহাব কোথা থেকে তার হুকুম নিয়েছে তা বোঝানো। সাহু সিজদা ও ভুল-সংশোধনের পূর্ণ রূপরেখা পার্ট ৫খ-এ।
\end{quote}

\end{document}
